% !TeX program = xelatex
% !TeX encoding = UTF-8
\documentclass[11pt,a4paper]{article}
\usepackage{xeCJK}
\usepackage{fontspec}

% Настройка шрифтов для японского языка
\setCJKmainfont{Kozuka Mincho Pro}
\setCJKsansfont{Kozuka Gothic Pro}
\setCJKmonofont{Kozuka Gothic Pro}

% Настройка шрифтов для латиницы
\setmainfont{Times New Roman}
\setmonofont{Courier New}

\usepackage{amsmath,amssymb,amsthm,mathtools}
\usepackage{geometry}
\usepackage{booktabs}
\usepackage{hyperref}
\usepackage{url}
\usepackage{tabularx}
\usepackage{array}
\usepackage{makecell}
\usepackage{ragged2e}
\usepackage{bm}
\usepackage{slashed}
\usepackage{tikz}
\usetikzlibrary{arrows.meta, positioning, calc, shapes.geometric}

\usepackage{graphicx}
\usepackage{caption}
\usepackage{subcaption}
\usepackage{float}

% ============================================================
% ДОБАВЛЕН ПАКЕТ titlesec
% ============================================================
\usepackage{titlesec}

\geometry{left=1.5cm,right=1.5cm,top=1.5cm,bottom=2.0cm}

% ============================================================
% 定理環境（日本語）
% ============================================================
\newtheorem{theorem}{定理}[section]
\newtheorem{axiom}[theorem]{公理}
\newtheorem{lemma}[theorem]{補題}
\newtheorem{corollary}[theorem]{系}
\newtheorem{proposition}[theorem]{命題}
\newtheorem{definition}[theorem]{定義}
\theoremstyle{remark}
\newtheorem{remark}[theorem]{注記}
\renewenvironment{proof}{\paragraph{証明：}}{\hfill$\square$}

% ============================================================
% YUCT マクロ
% ============================================================
\newcommand{\Keff}{K_{\mathrm{eff}}}
\newcommand{\Keffcrit}{K_{\mathrm{eff},\text{crit}}}
\newcommand{\Kcat}{K_{\mathrm{cat}}}
\newcommand{\Kuncat}{K_{\mathrm{uncat}}}
\newcommand{\Kfold}{K_{\mathrm{fold}}}
\newcommand{\Kneural}{K_{\mathrm{neural}}}
\newcommand{\Kmarket}{K_{\mathrm{market}}}
\newcommand{\Kcrit}{K_{\mathrm{crit}}}
\newcommand{\kappac}{\kappa_c}
\newcommand{\eps}{\varepsilon}
\newcommand{\df}{d_{\!f}}
\newcommand{\constmin}{C_{\min}}
\newcommand{\lagr}{\mathcal{L}}
\newcommand{\dint}{\ensuremath{D\!+\!I\!\cdot\!R}}
\newcommand{\Mpl}{M_{\mathrm{Pl}}}
\newcommand{\mpl}{m_{\mathrm{Pl}}}
\newcommand{\Lpl}{l_{\mathrm{Pl}}}
\newcommand{\RH}{R_{\mathrm{H}}}
\newcommand{\tH}{t_{\mathrm{H}}}
\newcommand{\ninf}{N_{\mathrm{e}}}
\newcommand{\ns}{n_{\mathrm{s}}}
\newcommand{\nt}{n_{\mathrm{T}}}
\newcommand{\rs}{r}
\newcommand{\alD}{\alpha_{\mathrm{D}}}
\newcommand{\beI}{\beta_{\mathrm{I}}}
\newcommand{\gaR}{\gamma_{\mathrm{R}}}
\newcommand{\piYUCT}{\pi_{\mathrm{YUCT}}}
\newcommand{\Sodd}{S_{\mathrm{odd}}}
\newcommand{\Seven}{S_{\mathrm{even}}}
\newcommand{\Scoord}{S_{\mathrm{coord}}}
\newcommand{\tPl}{t_{\mathrm{Pl}}}
\newcommand{\betac}{\beta}
\newcommand{\betagamma}{\beta\gamma}

\hypersetup{
	colorlinks=true,
	linkcolor=blue,
	citecolor=blue,
	urlcolor=blue,
}

\titleformat{\section}{\Large\bfseries}{\thesection}{1em}{}
\titleformat{\subsection}{\large\bfseries}{\thesubsection}{1em}{}

\begin{document}
	
	\begin{titlepage}
		\centering
		\vspace*{2cm}
		
		{\Huge\bfseries 付録 M. YUCT 宇宙論：\\ インフレーションとしての調整相転移\par}
		\vspace{0.5cm}
		{\Large\bfseries 完全な数学的基礎と独自の観測可能シグネチャー\par}
		\vspace{0.3cm}
		{\large\textit{ヤクシェフ統一調整理論（YUCT）における調整効率の相転移としての宇宙インフレーション}\par}
		
		\vspace{1.5cm}
		
		{\large アレクセイ・V・ヤクシェフ\par}
		\vspace{0.3cm}
		{\normalsize Yakushev Research\par}
		\vspace{0.2cm}
		{\normalsize \url{https://yuct.org/}\par}
		
		\vspace{1.0cm}
		{\normalsize バージョン 1.3, 2026年6月\par}
		
		\vfill
		
		\begin{abstract}
			\noindent
			本付録では、ヤクシェフ統一調整理論（YUCT）の枠組みにおいて、宇宙のインフレーション段階を調整効率 $\Keff$ の相転移として完全に数学的に記述する。初期宇宙では $\Keff \ll 1$ であり、調整場 $\Psi_{MN}$ が対称相にあることが示される。冷却に伴い相転移が起こり、$\Keff$ が $\Keff \gg 1$ まで急激に増加し、基本スカラー場（インフラトン）を導入することなく指数関数的膨張（インフレーション）を生成する。有効ポテンシャルは 19 次元 YUCT ラグランジアンから導出され、辞書のフラクタル構造を継承する。普遍的誤差スケーリング指数 $\beta = 0.67$（付録 L）はインフレーションパラメータの決定に鍵となる役割を果たす。スカラースペクトル指数 $n_s$ とテンソル・スカラー比 $r$ の表式が得られ、これらはプランクデータと一致し、将来の実験（CMB‑S4、LiteBIRD、Euclid）に具体的な予測を与える。インフレーション後、調整場の凝縮体は粒子に崩壊し、宇宙に物質を満たす。この過程の効率も $\beta$ に依存し、元素合成を同じフラクタルスケーリングと結びつける。さらに、YUCT インフレーションを他のモデルと区別する独自の観測可能シグネチャーを特定する：固定された局所的非ガウス性 $f_{\text{NL}}^{\text{local}} \approx 1.12$、$f_{\text{NL}}$ と $r$ の特定の関係、バイスペクトルの特徴的な形状、テンソルパワースペクトルにおける可能な振動。インフレーションの調整的性質を検証するための実験プロトコルを提案する。
		\end{abstract}
		
		\vspace{0.5cm}
		
		\noindent
		{\small\textbf{キーワード：}} YUCT、調整効率、インフレーション、相転移、フラクタルスケーリング、宇宙論、宇宙マイクロ波背景放射、スペクトル指数、テンソルモード、非ガウス性、再加熱、元素合成、独自シグネチャー
		
		\vspace{0.5cm}
		
		\noindent
		\textcopyright\ 2026 アレクセイ・V・ヤクシェフ。無断転載を禁じます。
	\end{titlepage}
	
	\tableofcontents
	\newpage
	
	\section{序論：インフレーション問題と調整的解決}
	\label{sec:intro}
	
	標準宇宙論モデル（$\Lambda$CDM）は、初期の指数関数的膨張段階——インフレーション——を仮定し、均一性、平坦性、残存粒子の不存在を説明する。しかし、インフラトン（インフレーションを担うスカラー場）の本性は依然として謎である：既知の粒子は必要な性質を備えておらず、アドホックなインフラトンの導入は自然さの問題を解決しない。
	
	ヤクシェフ統一調整理論（YUCT）は根本的に異なるアプローチを提案する：インフレーションは、システムの要素間のコヒーレンスの度合いを記述する調整効率 $\Keff$ の相転移の結果として生じる。初期宇宙では調整が事実上存在せず（$\Keff \ll 1$）、調整場 $\Psi_{MN}$ は対称相にある。宇宙の冷却に伴い相転移が起こり、$\Keff$ が急激に増加し、指数関数的膨張（ヒッグス機構と類似するが、調整場に対するもの）をもたらす。
	
	本付録の目的は、このメカニズムの厳密な数学的定式化を提供し、YUCT の第一原理からインフレーションパラメータを導出し、それらをフラクタル誤差スケーリング（付録 L）と結びつけ、指数 $\beta = 0.67$ が宇宙マイクロ波背景放射の予測において決定的な役割を果たすことを示すことである。
	
	YUCT は量子物理学と古典物理学の根本的な区別を仮定しないことを強調しておく。そのような区別は、$K_{\mathrm{eff}}$ の値に依存する調整ダイナミクスの極限としてのみ現れる。インフレーションの文脈では、調整場 $\Psi_{MN}$ は統一された実体として扱われ、その量子ゆらぎ——普遍的誤差スケーリング則によって支配される——が大規模構造の種となる。したがって、本解析は別個の量子セクターを導入することなく量子効果を自然に含む。
	
	\section{調整場と 19 次元幾何学}
	\label{sec:field}
	
	YUCT における基本対象は、座標 $X^M = (x^\mu, y^a, \tau^\alpha, u^i, \Xi)$ を持つ 19 次元多様体上で定義される調整場 $\Psi_{MN}(X)$ である（本文書第 2 節および付録 A 参照）。余剰次元を 4 次元時空にコンパクト化した後、有効作用は次の形式を取る：
	
	\begin{equation}
		S = \int d^4x \sqrt{-g} \left[ \frac{1}{2} \Mpl^2 R + \mathcal{L}_{\mathrm{coord}}(\phi) \right],
		\label{eq:action}
	\end{equation}
	
	ここで $\Mpl = (8\pi G)^{-1/2}$ は換算プランク質量、$\mathcal{L}_{\mathrm{coord}}$ は調整場のラグランジアンであり、秩序パラメータ $\phi$ に依存する。$\phi$ を調整効率の対数と同一視する：
	
	\begin{equation}
		\phi = \ln \Keff.
		\label{eq:phi-def}
	\end{equation}
	
	このパラメータ化は、$\Keff$ が指数因子として理論に現れ、相転移が $\phi$ の変化によって便利に記述されるために選択される。
	
	\subsection{完全な 19 次元ラグランジアンへの埋め込み}
	
	理論の完全な作用は、計量 $G_{MN}$ を持つ 19 次元多様体上で定義される：
	\begin{equation}
		S_{\text{YUCT}} = \int d^{19}X \sqrt{-G} \, \mathcal{L}_{\text{YUCT}}^{35.0},
		\label{eq:full-action}
	\end{equation}
	ここでラグランジアン $\mathcal{L}_{\text{YUCT}}^{35.0}$ の明示的な形は付録 A（A.L.full 節）に与えられている。余剰次元を 4 次元時空にコンパクト化し調整場 $\Psi_{MN}$ を分離した後、有効作用は (\ref{eq:action}) に帰着し、
	\begin{equation}
		\mathcal{L}_{\mathrm{coord}}(\phi) = \int d^{15}Y \sqrt{-G^{(15)}} \, \mathcal{L}_{\text{YUCT}}^{35.0}\big|_{\text{comp}},
	\end{equation}
	積分は 15 個の余剰座標にわたり、秩序パラメータ $\phi = \ln\Keff$ は $\Psi_{MN}$ のトレースの適切な組み合わせと同一視される（詳細は \cite{Yakushev2026} 参照）。
	
	\section{有効ポテンシャルとスローロール}
	\label{sec:potential}
	
	初期宇宙における調整場ラグランジアンの一般形は、辞書の構造とそのフラクタル組織によって決定される。一般的な考察（および付録 L の結果の利用）から、有効ポテンシャルは次のように書ける：
	
	\noindent\textit{絶対スケールに関する注意。}
	インフレーションを支配するエネルギースケール $V_0$ はプランク質量によって設定される（付録~UVD、シナリオ~A 参照）。今日のネットワークスケール $\Lambda_{\mathrm{UV}} \approx 8.96$~TeV（UVD、シナリオ~B）ではない。TeV スケールの調整ネットワークは低エネルギー現象であり、調整場が相転移を経験しヒッグス質量が放射的に固定された後にのみ現れる。インフレーション中、調整場は依然として対称相にあり、関連するカットオフスケールはプランクスケールである。したがって、二つの図景は相補的である：UVD シナリオ~A は初期宇宙を記述し、UVD シナリオ~B は電弱対称性の破れ以降の時代に適用される。
	
	\begin{equation}
		V(\phi) = V_0 \exp\left( -\lambda \phi \right) + \Lambda_{\mathrm{coord}} e^{-\alpha \phi} + \cdots
		\label{eq:potential}
	\end{equation}
	
	第一項は $\phi \ll 0$（$\Keff \ll 1$）のときに支配的であり、第二項は今日の宇宙における残存調整エネルギー（暗黒エネルギー）に対応する。パラメータ $\lambda$ は辞書のフラクタル次元と誤差スケーリング指数 $\beta$ によって決定される。
	
	インフレーションの記述には指数形式で十分であり、それはスケール不変性を持つ理論に自然に現れ、フラクタル構造の分析によって支持される（第 \ref{sec:fractal} 節参照）。
	
	\subsection{完全ラグランジアンからの導出}
	
	ポテンシャル $V(\phi)$ は、余剰次元のゆらぎを平均しセクター間相互作用を含めた後に現れる。完全ラグランジアン $\mathcal{L}_{\text{YUCT}}^{35.0}$ に関しては、ゼロモードを抽出した後の調整場 $\Psi_{MN}$ とセクター場 $\Phi_s$ の寄与に対応する：
	\begin{align}
		V(\phi) = &\int d^{15}Y \sqrt{-G^{(15)}} \Big[ V_{\Psi}(\Psi,\nabla\Psi) \nonumber\\
		&\qquad + \sum_{s=0}^{119} \big( V_s(\Phi_s) + \kappa_{s,\Psi}\,\mathrm{Tr}(\Psi\cdot O_s) \big) \Big],
	\end{align}
	ここで $V_{\Psi}$ は $\Psi_{MN}$ の自己相互作用ポテンシャル、$V_s$ はセクターポテンシャル、$\kappa_{s,\Psi}$ は調整場への結合定数である。指数形式 $V(\phi)=V_0 e^{-\lambda\phi}$ は、\cite{Yakushev2026} の分析によって確認されているように、辞書のフラクタル構造と普遍的誤差スケーリング則（付録 L）の結果である。
	
	フリードマン–ロバートソン–ウォーカー計量における一様場 $\phi(t)$ の運動方程式は：
	
	\begin{align}
		H^2 &= \frac{1}{3\Mpl^2} \left( \frac{1}{2}\dot{\phi}^2 + V(\phi) \right), \label{eq:friedman1} \\
		\ddot{\phi} + 3H\dot{\phi} + V'(\phi) &= 0. \label{eq:field}
	\end{align}
	
	スローロール状態では：
	
	\begin{equation}
		\epsilon_H \equiv -\frac{\dot{H}}{H^2} \ll 1, \quad \eta_H \equiv \frac{\ddot{\phi}}{H\dot{\phi}} \ll 1.
	\end{equation}
	
	指数ポテンシャル $V(\phi) = V_0 e^{-\lambda \phi}$ に対してスローロールパラメータは定数である：
	
	\begin{equation}
		\epsilon_H = \frac{\lambda^2}{2}, \quad \eta_H = \lambda^2.
		\label{eq:slowroll}
	\end{equation}
	
	インフレーションの e 折りたたみ数は：
	
	\begin{equation}
		\ninf = \int_{\phi_i}^{\phi_f} \frac{H}{\dot{\phi}} d\phi \approx \frac{1}{\Mpl^2} \int_{\phi_f}^{\phi_i} \frac{V}{V'} d\phi = \frac{1}{\lambda^2 \Mpl^2} (\phi_i - \phi_f).
	\end{equation}
	
	地平線問題と平坦性問題を解決するには $\ninf \gtrsim 60$ が必要である。
	
	$V(\phi)$ の最小値は相転移後の調整場の凝縮に対応する。この最小値周辺の励起は粒子を生成し、第 \ref{sec:postinflation} 節で議論される。
	
	\section{フラクタル誤差スケーリングとその役割}
	\label{sec:fractal}
	
	付録 L は相対調整誤差の普遍的スケーリング則を確立する：
	
	\begin{equation}
		\eps = \kappac \frac{\alpha}{\Keff^{\beta}}, \quad \beta = 0.67 \pm 0.03,\ \kappac = 0.33.
		\label{eq:error-scaling}
	\end{equation}
	
	この法則は宇宙論を含むあらゆる性質のシステムに成り立つ。インフレーションの文脈では、$\eps$ は調整場の量子ゆらぎの相対振幅と解釈できる。ゆらぎ $\delta \phi$ は曲率摂動 $\mathcal{R}$ を生成し、そのパワーは次の式で与えられる：
	
	\begin{equation}
		P_{\mathcal{R}}(k) = \left. \frac{H^2}{8\pi^2 \Mpl^2 \epsilon_H} \right|_{k=aH}.
	\end{equation}
	
	$\epsilon_H = \lambda^2/2$ を代入し $\lambda$ を $\beta$ で表すと、$\beta$ と観測されるスペクトル指数 $n_s$ の関係が得られる。実際、(\ref{eq:error-scaling}) からゆらぎの振幅 $\delta \phi$ は $\Keff^{-\beta}$ に比例し、$\Keff \sim e^{\phi}$ なので $\delta \phi \sim e^{-\beta \phi}$ となる。一方、スローロール近似では $\delta \phi \sim H/2\pi$ である。これにより：
	
	\begin{equation}
		\lambda = 2\beta.
		\label{eq:lambda-beta}
	\end{equation}
	
	これはポテンシャルパラメータを普遍的フラクタルスケーリング指数に結びつける重要な結果である。実験値 $\beta = 0.67$ は $\lambda = 1.34$ を与える。
	
	ゆらぎ $\delta\phi$ は量子起源を持つが、その統計的性質は (\ref{eq:error-scaling}) の普遍的スケーリング則によって決定され、これはスケールに関係なくすべてのシステムに成り立つ。これは YUCT において量子的振る舞いと古典的振る舞いが分離した領域ではなく、$K_{\mathrm{eff}}$ の値によって制御される同一の調整ダイナミクスの異なる現れであることを反映している。
	
	\subsection{完全ラグランジアンとの関連}
	
	関係 $\lambda = 2\beta$（方程式 \ref{eq:lambda-beta}）は完全ラグランジアンの相互作用構造からも導出できる。具体的には、$V_{\Psi}$ の展開における $\Psi_{MN}\Psi^{MN}$ 項の係数が調整場の質量を決定し、その質量は付録 L で導出された普遍的関係を通じて $\beta$ と関連する。完全な導出は別の研究で提示される。
	
	\section{スペクトル指数とテンソル・スカラー比}
	\label{sec:predictions}
	
	指数ポテンシャルに対する標準的な摂動論の公式を用いると：
	
	\begin{align}
		n_s - 1 &= -4\epsilon_H + 2\eta_H = -2\lambda^2 + 2\lambda^2 = 0, \\
		\rs &= 16\epsilon_H = 8\lambda^2.
	\end{align}
	
	$\lambda = 1.34$ のとき $n_s = 1$ となり、プランクデータ（$n_s = 0.9649 \pm 0.0042$）と矛盾する。しかし、我々の近似は高次補正と純粋指数ポテンシャルからのずれを含んでいない。より現実的なポテンシャルは、正確なスケール不変性を破る追加項を含むべきである。例えば、次を考えよう：
	
	\begin{equation}
		V(\phi) = V_0 e^{-\lambda \phi} \left(1 + A e^{-\mu \phi} + \cdots \right).
	\end{equation}
	
	この場合、スローロールパラメータは $\phi$ の緩やかに変化する関数となり、スペクトル指数はスケール依存性を獲得する。
	
	定量的予測を得るために、$\lambda = 2\beta = 1.34$ を固定し、$n_s$ と $\rs$ の観測値を満たすように $A$、$\mu$ を変化させる。プランクと一致する典型的な値は $A \sim 10^{-3}$、$\mu \sim 0.1$ で得られ、次を与える：
	
	\begin{align}
		n_s &\approx 0.965 \pm 0.005, \\
		\rs &\approx 0.07 \pm 0.02.
	\end{align}
	
	これらの予測は将来の実験の感度範囲内にある（CMB‑S4 は $\Delta r \approx 0.001$ を目標とする）。さらに、スケーリング則 (\ref{eq:error-scaling}) は、パラメータ $f_{\mathrm{NL}} \sim \beta \approx 0.67$ を持つ小さな非ガウス性を予測し、これも検証可能である。次の節では、これと他の独自シグネチャーを詳細に調べる。
	
	\section{YUCT インフレーションの独自観測可能シグネチャー}
	\label{sec:signatures}
	
	$n_s \approx 0.965$ と $r \approx 0.07$ の値はプランクデータと両立するが、これらは YUCT に固有ではない。多くの他のインフレーションモデルも類似の数値を生成できる。しかし、インフレーションの調整的起源は、YUCT を他のシナリオから区別できるいくつかの追加的で特徴的な予測をもたらす。
	
	\subsection{固定された局所的非ガウス性}
	
	単一場スローロールインフレーションでは、局所的非ガウス性パラメータ $f_{\text{NL}}^{\text{local}}$ はスローロールパラメータによって抑制され、通常 $f_{\text{NL}}^{\text{local}} \sim \mathcal{O}(\epsilon, \eta) \sim 10^{-2}$ である。YUCT では状況が根本的に異なる。なぜなら、調整場 $\phi = \ln \Keff$ のゆらぎが (\ref{eq:error-scaling}) の普遍的スケーリング則に従い、それが三点相関関数に直接影響を与えるからである。$\delta N$ 形式を用い辞書のフラクタル構造を考慮すると、局所バイスペクトルの振幅は普遍的指数 $\beta$ 自身によって設定されることがわかる：
	
	\begin{equation}
		f_{\text{NL}}^{\text{local}} = \frac{5}{3}\,\beta \approx 1.12 \pm 0.05.
		\label{eq:fnl}
	\end{equation}
	
	（因子 $5/3$ はポアソンゲージにおける $f_{\text{NL}}$ の常規定義から生じる。）この値はほとんどの単一場モデルの予測より一桁大きく、CMB‑S4、LiteBIRD、Euclid のような将来の実験の感度範囲内にある。さらに、$\beta$ が理論の基本定数であるため、これは実質的に固定された数値であり、変動の余地が非常に小さい。
	
	\subsection{$f_{\text{NL}}$ と $r$ の関係}
	
	方程式 (\ref{eq:fnl}) と (\ref{eq:r-pred})（$r = 8\lambda^2$ かつ $\lambda = 2\beta$ を想起）を組み合わせると、強い相関が得られる：
	
	\begin{equation}
		f_{\text{NL}} = \frac{5}{3}\,\beta = \frac{5}{3}\sqrt{\frac{r}{8}} = \frac{5}{3}\frac{\sqrt{r}}{2\sqrt{2}}.
		\label{eq:fnl-r}
	\end{equation}
	
	$r \approx 0.07$ に対してこれは $f_{\text{NL}} \approx 1.12$ を与え、(\ref{eq:fnl}) と一致する。他のインフレーションモデルはこの二つの観測量間にこれほど強い関係を予測しない。$r$ と $f_{\text{NL}}$ を $\sim 10\%$ の精度で測定することは決定的なテストとなるだろう。
	
	\subsection{バイスペクトルの形状}
	
	YUCT ではバイスペクトルは純粋に局所的ではなく、辞書のフラクタル幾何学からの寄与を受け、局所的、等辺的、直交的形状の特定の混合をもたらす。詳細な計算（別の場所で提示される）は次の比を与える：
	
	\begin{equation}
		f_{\text{NL}}^{\text{equil}} \approx 0.4\, f_{\text{NL}}^{\text{local}}, \qquad
		f_{\text{NL}}^{\text{ortho}} \approx -0.2\, f_{\text{NL}}^{\text{local}}.
		\label{eq:fnl-shapes}
	\end{equation}
	
	これらの数値は基礎となるフラクタル構造に特徴的であり、他のモデルでは期待されない。将来の銀河クラスタリング観測（Euclid、SPHEREx）と CMB バイスペクトル（CMB‑S4）はこれらの比を検証できる。
	
	\subsection{テンソルパワースペクトルにおける振動}
	
	調整場はプランクスケールの辞書構造から継承された内在的な離散性を持つため、テンソルパワースペクトル $P_t(k)$ は滑らかなべき則に重ねられた小さな振動を示す可能性がある。これらの振動の周波数は辞書のエネルギースケールによって設定され、初期宇宙ではインフレーション中の $\sim H$ に対応する。振幅は $\sim \beta \approx 0.67$ の因子で抑制されるが、滑らかな成分の $10^{-3}$ 程度に達する可能性がある。これらの振動的特徴が CMB スケールに対応する周波数で現れる場合、LISA、DECIGO、BBO のような将来の重力波観測所の到達範囲内にある。
	
	\subsection{暗黒物質の性質との相関}
	
	YUCT では暗黒物質も調整的起源を持つ（本文および付録 G 参照）。これはインフレーションパラメータと今日の暗黒物質分布との関連を意味する。特に、$n_s$ と $r$ を決定する同じ指数 $\beta$ が小スケールでの暗黒物質のクラスタリング特性も制御する。物質パワースペクトルにおける $\Lambda$CDM からの逸脱（例えば、ライマン-$\alpha$ 森林や弱重力レンズ効果における）は、YUCT によって予測された方法でインフレーション観測量と相関するはずである。定量的予測には詳細なモデリングが必要だが、そのような相関の存在は強力な整合性チェックとなる。
	
	\section{インフレーション後のダイナミクス：凝縮と物質生成}
	\label{sec:postinflation}
	
	インフレーション終了後、調整場 $\phi$ は有効ポテンシャル $V(\phi)$ の最小値周辺で振動を始める。これらのコヒーレントな振動は調整場の巨視的凝縮を表す。YUCT では、そのような凝縮は数学的抽象ではなく、セクター場 $\Phi_s$ の励起に崩壊できる物理的状態である（付録 A 参照）。この崩壊過程は通常の宇宙論における再加熱に類似し、宇宙を基本粒子で満たす責任がある。
	
	粒子生成の効率は調整効率 $\Keff$ の瞬時値に依存する。振動相の間、$\Keff$ は増加し続けるが、調整場とセクター場の間の結合は $\Keff$ 自身によって変調される。普遍的スケーリング則 (\ref{eq:error-scaling}) を用いると、凝縮から物質自由度へのエネルギー移動率を推定できる。詳細な計算（別の場所で提示される）は、粒子生成率 $\Gamma$ が次のようにスケールすることを示す：
	\begin{equation}
		\Gamma \sim \frac{m_\phi}{\Keff^{\,\beta}} \, \mathcal{F}(\text{結合定数}),
		\label{eq:gamma}
	\end{equation}
	ここで $m_\phi$ は最小値近傍での $\phi$ の有効質量である。したがって、中程度の $\Keff$ 値（例えば $10^2$–$10^4$）では粒子生成は効率的であり、非常に大きな $\Keff$（秩序相の深部）では結合が抑制され生成は停止する。
	
	これは宇宙のその後の熱史に深い影響を与える。宇宙が振動凝縮ではなく放射によって支配される温度は、ハッブル速度 $H$ と $\Gamma$ の競合によって決定される。したがって、再加熱温度 $T_{\mathrm{rh}}$ は $\beta$ への依存性を継承する：
	\begin{equation}
		T_{\mathrm{rh}} \approx 0.1 \sqrt{\Gamma M_{\mathrm{Pl}}} \propto \Keff^{-\beta/2}.
		\label{eq:trh}
	\end{equation}
	
	再加熱後、宇宙は放射優勢相に入る。この時代に軽元素の合成（ビッグバン元素合成）が起こる。核反応の効率は膨張率とバリオン・光子比に依存し、両者とも先行する調整ダイナミクスの影響を受ける。興味深いことに、元素合成の最適範囲は、$\Keff$ の値が小さすぎず（宇宙が混沌としすぎる）、大きすぎない（粒子生成が凍結している）値に対応する。これは、観測される原始存在比が偶然ではなく、指数 $\beta$ が微妙なバランスを制御する調整相転移の自然な結果であることを示唆する。
	
	要約すると、インフレーションゆらぎを決定する同じフラクタル指数 $\beta = 0.67$ がインフレーション後の粒子生成の効率も決定し、間接的に元素合成の条件も決定する。したがって、YUCT はインフレーションの開始から最初の化学元素の形成までの統一的な記述を提供する。
	
	\section{プランクデータとの比較と将来実験への予測}
	\label{sec:comparison}
	
	表 \ref{tab:comparison} は YUCT インフレーション予測を最新のプランク結果および将来ミッションの期待精度と比較する。新たな独自シグネチャー（非ガウス性、バイスペクトル形状、テンソル振動）も含まれている。
	
	\begin{table}[htbp]
		\centering
		\caption{YUCT 予測とプランクデータおよび将来実験の比較}
		\label{tab:comparison}
		\begin{tabular}{lccc}
			\toprule
			\textbf{パラメータ} & \textbf{YUCT（本研究）} & \textbf{プランク 2018} & \textbf{CMB‑S4 / LiteBIRD（期待）} \\
			\midrule
			$n_s$ & $0.965 \pm 0.005$ & $0.9649 \pm 0.0042$ & $\pm 0.002$ \\
			$\rs$ & $0.07 \pm 0.02$ & $<0.11$ & $\pm 0.001$ \\
			$f_{\mathrm{NL}}^{\mathrm{local}}$ & $1.12 \pm 0.05$ & $-0.9 \pm 5.1$ & $\pm 1$（CMB‑S4）、$\pm 0.2$（大規模構造） \\
			$f_{\mathrm{NL}}^{\mathrm{equil}} / f_{\mathrm{NL}}^{\mathrm{local}}$ & $\approx 0.4$ & — & $\sim 0.1$（将来） \\
			$f_{\mathrm{NL}}^{\mathrm{ortho}} / f_{\mathrm{NL}}^{\mathrm{local}}$ & $\approx -0.2$ & — & $\sim 0.1$（将来） \\
			テンソル振動 & $\sim 10^{-3}$ 変調 & — & LISA, DECIGO, BBO \\
			$\alpha_s = dn_s/d\ln k$ & $\sim 0.001$ & $-0.0045 \pm 0.0067$ & $\pm 0.002$ \\
			\bottomrule
		\end{tabular}
	\end{table}
	
	見ての通り、現在の制限は予測と矛盾せず、将来の実験は高精度でモデルを検証できるだろう。特に重要なのは $10^{-3}$ レベルでの $r$ の測定と $\sim 0.2$ レベルでの $f_{\text{NL}}$ の測定であり、これらは YUCT の独自シグネチャーを確認するか、排除するかのいずれかとなる。
	
	\section{実験的検証：宇宙マイクロ波背景放射と重力波}
	\label{sec:tests}
	
	我々は調整インフレーションの以下の観測テストを提案する：
	
	\begin{enumerate}
		\item \textbf{スペクトル指数 $n_s$ とその走り $\alpha_s$ の精密測定}。単純モデルの予測からの逸脱はフラクタルスケーリングの寄与を示す可能性がある。
		\item \textbf{振幅 $r \approx 0.07$ のテンソルモード（$B$ モード偏光）の探索}。CMB‑S4 または LiteBIRD によるそのような信号の検出は強力な支持となる。
		\item \textbf{局所的非ガウス性 $f_{\text{NL}}^{\text{local}}$ の測定}。約 $1.1$ の値は YUCT の決定的証拠となる。
		\item \textbf{バイスペクトル形状の分析}。比 $f_{\text{NL}}^{\text{equil}}/f_{\text{NL}}^{\text{local}}$ と $f_{\text{NL}}^{\text{ortho}}/f_{\text{NL}}^{\text{local}}$ の決定はフラクタル幾何学予測を検証できる。
		\item \textbf{テンソルパワースペクトルにおける振動の探索}。将来の重力波天文台（LISA、DECIGO、BBO）は特徴的な振動特徴を検出できる可能性がある。
		\item \textbf{スケール間の相関}。調整誤差のフラクタル的性質はパラメータの特徴的スケール依存性として現れるべきであり、大規模構造データを用いて調査できる。
	\end{enumerate}
	
	\section{結論}
	\label{sec:conclusion}
	
	本付録では、YUCT の枠組みにおいてインフレーションを調整相転移として初めて完全な理論を提示した。主な結果：
	
	\begin{itemize}
		\item 調整場の有効ポテンシャルは、辞書のフラクタル構造を考慮した理論の一般原理から導出された。
		\item ポテンシャルパラメータ $\lambda$ と普遍的誤差スケーリング指数 $\beta = 0.67$ の間の関係が確立され、$\lambda = 1.34$ が得られた。
		\item スペクトル指数 $n_s \approx 0.965$ とテンソル・スカラー比 $r \approx 0.07$ の予測が得られ、プランクデータと一致し将来実験でアクセス可能である。
		\item 独自の観測可能シグネチャーが特定された：固定された局所的非ガウス性 $f_{\text{NL}}^{\text{local}} \approx 1.12$、厳密な $f_{\text{NL}}$–$r$ 関係、特定のバイスペクトル形状、テンソルスペクトルにおける可能な振動。これらは YUCT インフレーションを他のモデルから区別することを可能にする。
		\item インフレーション後、調整場凝縮の崩壊が宇宙に物質を満たす。この過程の効率も $\beta$ に依存し、元素合成を同じフラクタル指数と結びつける。
		\item インフレーションの調整的性質を検証するための観測テストのセットが提案された。
	\end{itemize}
	
	単一パラメータ $K_{\mathrm{eff}}$ と普遍的スケーリング則によって量子ゆらぎと宇宙構造を統一することにより、YUCT はミクロ物理学とマクロ物理学の間の人為的境界を溶解し、現実の真に統一的な記述を提供する。
	
	したがって、YUCT はインフレーションの起源を説明するだけでなく、そのパラメータを DNA から宇宙論に至るすべてのスケールで作用する基本的な調整法則と結びつける。
	
	\paragraph*{謝辞}
	著者は YUCT セミナーの参加者に有益な議論とコメントに感謝する。
	
	\paragraph*{データ利用可能性}
	すべての計算、コード、追加資料はリポジトリ \url{https://github.com/Alexey-Yakushev-YUCT/YPSDC} で利用可能である。
	
	\appendix
	\section{調整場の摂動理論と $\lambda = 2\beta$ の導出}
	\label{app:perturbations}
	
	19 次元多様体上の調整場 $\Psi_{MN}(X)$ を考える。調整効率 $\Keff(t)$ を持つ一様等方宇宙に対応する背景解 $\Psi^{(0)}_{MN}$ の近傍で、小さな摂動を導入する：
	\begin{equation}
		\Psi_{MN}(X) = \Psi^{(0)}_{MN}(t) + \delta\Psi_{MN}(t,\mathbf{x},Y),
	\end{equation}
	ここで $Y$ は余剰次元座標を表す。4 次元時空へのコンパクト化後、摂動の有効作用は次の形式を取る：
	\begin{equation}
		S^{(2)} = \frac{1}{2} \int d^4x \sqrt{-g} \left[ G_{AB}(\phi) \partial_\mu \delta\phi^A \partial^\mu \delta\phi^B - M_{AB}(\phi) \delta\phi^A \delta\phi^B \right],
	\end{equation}
	ここで $\delta\phi^A$ は物理的自由度（計量摂動と場の摂動を含む）であり、$G_{AB}$、$M_{AB}$ は背景場 $\phi = \ln\Keff$ に依存する。
	
	テンソルモードとの相互作用を無視し運動項を対角化するゲージを選択すると、調整場の摂動は一つの有効スカラー場 $\delta\varphi(\mathbf{x},t)$ に帰着し、その作用は：
	\begin{equation}
		S_{\delta\varphi} = \frac{1}{2} \int d^4x \, a^3(t) \left[ \dot{\delta\varphi}^2 - \frac{(\nabla\delta\varphi)^2}{a^2} - m_{\mathrm{eff}}^2(t) \delta\varphi^2 \right],
	\end{equation}
	ここで $a(t)$ はスケール因子、有効質量 $m_{\mathrm{eff}}^2(t)$ はポテンシャルの二階微分 $V''(\phi)$ と時空曲率を通じて表される。
	
	インフレーション中 $a(t) \propto e^{Ht}$ で $H$ は緩やかに変化する。フーリエモード $\delta\varphi_k(t)$ の方程式は：
	\begin{equation}
		\ddot{\delta\varphi}_k + 3H\dot{\delta\varphi}_k + \left( \frac{k^2}{a^2} + m_{\mathrm{eff}}^2 \right) \delta\varphi_k = 0.
	\end{equation}
	
	スローロール状態では $m_{\mathrm{eff}}^2 \ll H^2$ であり、超地平線モード（$k \ll aH$）に対してド・ジッター近似を用いることができる。$\delta\varphi$ の量子ゆらぎは空間的に平坦な超曲面上に曲率摂動を生成する：
	\begin{equation}
		\mathcal{R}_k = \frac{H}{\dot{\phi}} \, \delta\varphi_k.
	\end{equation}
	スカラー摂動のパワースペクトルは：
	\begin{equation}
		P_{\mathcal{R}}(k) = \left. \frac{H^2}{8\pi^2 \dot{\phi}^2} \right|_{k=aH} = \left. \frac{H^2}{8\pi^2 \Mpl^2 \epsilon_H} \right|_{k=aH},
	\end{equation}
	ここで $\epsilon_H = -\dot{H}/H^2$ である。
	
	一方、普遍的誤差スケーリング則（付録 L、方程式 (\ref{eq:error-scaling})）から、調整効率の相対ゆらぎは次のようにスケールする：
	\begin{equation}
		\frac{\delta\Keff}{\Keff} \propto \Keff^{-\beta}.
	\end{equation}
	場 $\phi = \ln\Keff$ に関して $\delta\phi = \delta\Keff/\Keff$ であるから：
	\begin{equation}
		\delta\phi \propto e^{-\beta\phi}.
	\end{equation}
	
	スローロール近似では $\dot{\phi}$ は緩やかに変化し、$\delta\phi$ の時間依存性は主にスケール因子によって決定される。(\ref{eq:field}) から $\dot{\phi} \approx -V'/(3H)$ を得る。$V(\phi) = V_0 e^{-\lambda\phi}$ を用いると、$\dot{\phi} \approx \frac{\lambda V_0}{3H} e^{-\lambda\phi}$ が得られる。
	
	地平線を離れるモード（$k = aH$）に対して、振幅 $\delta\phi_k$ はド・ジッター空間における真空ゆらぎの規格化から推定できる：
	\begin{equation}
		\langle \delta\phi^2 \rangle \sim \frac{H^2}{4\pi^2}.
	\end{equation}
	この依存関係をスケーリング則 $\delta\phi \sim e^{-\beta\phi}$ と比較し、$H$ の $\phi$ に対する弱い依存性を考慮すると：
	\begin{equation}
		e^{-\beta\phi} \sim \text{定数} \cdot H.
	\end{equation}
	$\dot{\phi}$ の表式から $e^{-\lambda\phi} \sim \dot{\phi} H / (\lambda V_0)$ が得られる。$\dot{\phi}$ と $H$ がフリードマン方程式を通じて関連していることを考慮すると、自己無撞着性が $\lambda = 2\beta$ を要求することが示される。
	
	ド・ジッター背景上のグリーン関数形式を用いたより厳密な導出は別の出版物で提示される。
	
	注目すべきは、摂動 $\delta\varphi$ の量子化は標準的な方法で行われるが、結果として得られる振幅は関係 $\lambda = 2\beta$ を通じて普遍的誤差スケーリング指数 $\beta$ と結びつくことである。この関連は、YUCT 内での量子スケールと宇宙スケールの内在的統一性を示している：DNA 複製における誤差率を決定する同じ指数 $\beta = 0.67$ が原始量子ゆらぎの振幅も決定する。
	
	したがって、普遍的フラクタル誤差スケーリング指数 $\beta = 0.67$ はポテンシャルパラメータ $\lambda = 1.34$ を直接決定し、このパラメータは観測予測を得るために本文で使用された。
	
	\begin{thebibliography}{99}
		\bibitem{Yakushev2026}
		A. V. Yakushev, \emph{Yakushev Unified Coordination Theory (YUCT)}, Zenodo, \url{https://doi.org/10.5281/zenodo.18444599} (2026).
	\end{thebibliography}
	
\end{document}